\section{Arrow of time under observation}\label{sec:arrow-protocol}

This section proves \thmNGArrowDPI{} and \thmNGProtocolTrap.
Both rest on two standard facts: relative entropy cannot increase under data processing \citep{CoverThomas2006}, and a chain in detailed balance looks the same forwards and backwards \citep{Kelly1979}.
For a stationary chain the arrow is the entropy production over the observation window \citep{Seifert2012,MaesNetocny2003}, and the results then say that coarse observation can hide dissipation but never create it \citep{Esposito2012}.

\subsection{Observation cannot create an arrow}

\begin{lemma}[Deterministic data processing]\label{lem:deterministic_dpi_kl}
Let $g:\Omega\to\Omega'$ be a map between finite sets. For all $p,q\in\Delta(\Omega)$,
\[
  \DKL{g_{\#}p}{g_{\#}q}\le\DKL{p}{q}\qquad\text{in }[0,\infty].
\]
\end{lemma}

\begin{proof}
See \Cref{app:probkl}; the proof applies the log-sum inequality fiber by fiber.
\end{proof}

\begin{lemma}[Reversal commutes with observation]\label{lem:rev_pushforward_commute}
For every $T\ge 0$, $f^{(T)}\circ\rev=\rev\circ f^{(T)}$ as maps $X^{T+1}\to Y^{T+1}$. Consequently, for every $Q\in\Delta(X^{T+1})$,
\[
  f^{(T)}_{\#}\bigl(\rev_{\#}Q\bigr)=\rev_{\#}\bigl(f^{(T)}_{\#}Q\bigr).
\]
\end{lemma}

\begin{proof}
Both composites send $(x_0,\dots,x_T)$ to $(f(x_T),\dots,f(x_0))$.
Pushforward respects composition, so $f^{(T)}_{\#}\rev_{\#}Q=(f^{(T)}\circ\rev)_{\#}Q=(\rev\circ f^{(T)})_{\#}Q=\rev_{\#}f^{(T)}_{\#}Q$.
\end{proof}

\begin{proof}[Proof of \Cref{thm:NG_ARROW_DPI}]
Write $Q=\mathbf{P}^{\mu,P}_T$.
By \Cref{lem:rev_pushforward_commute}, the observed path law and its reversal are the images of $Q$ and $\rev_{\#}Q$ under the single map $g=f^{(T)}$:
\[
  \mathbf{P}^{\mu,P,f}_T=g_{\#}Q,\qquad \rev_{\#}\mathbf{P}^{\mu,P,f}_T=g_{\#}\bigl(\rev_{\#}Q\bigr).
\]
\Cref{lem:deterministic_dpi_kl} with $p=Q$ and $q=\rev_{\#}Q$ gives $\Arr_T(\mu,P;f)\le\Arr_T(\mu,P)$.
\end{proof}

The only property of the observation used in the proof is that it commutes with reversal.
This identifies which observations are covered and which are not.

\begin{remark}[Noisy observation is covered; time-dependent observation is not]\label{rem:arrow_scope}
Suppose each $X_t$ is observed through the same noisy channel $C$, a Markov kernel from $X$ to $Y$, independently at each time.
The observed path law is then the image of $\mathbf{P}^{\mu,P}_T$ under the product channel $C^{\otimes(T+1)}$.
That channel commutes with reversal, and the data-processing inequality holds for arbitrary channels \citep{CoverThomas2006}, so the conclusion of \thmNGArrowDPI{} still holds.
Leaving the deterministic setting is therefore not by itself a way to create an arrow.

An observation that depends on time is different.
Let $X$ have a single state, so that every arrow of the chain is zero, and record $0$ at time $0$ and $1$ at time $1$.
The observed path is $(0,1)$ with probability one and its reversal $(1,0)$ has probability zero, so the observed arrow is infinite.
Here the arrow comes from a clock in the observer, not from the system.
A Markov chain fitted to the observed process is a third kind of object; the theorem makes no claim about it.
\end{remark}

\subsection{The protocol trap}

\begin{lemma}[Detailed balance gives reversible paths]\label{lem:detailed_balance_path}
If $\pi$ satisfies detailed balance with $P$, then for every $T\ge 0$ and every path $(x_0,\dots,x_T)$,
\[
  \pi(x_0)\prod_{t=0}^{T-1}P(x_t,x_{t+1})=\pi(x_T)\prod_{t=0}^{T-1}P(x_{t+1},x_t).
\]
Equivalently, $\mathbf{P}^{\pi,P}_T=\rev_{\#}\mathbf{P}^{\pi,P}_T$.
\end{lemma}

\begin{proof}
See \Cref{app:probkl}.
The proof is an induction on $T$ that uses detailed balance once per step and never divides, so it covers states with $\pi(x)=0$.
\end{proof}

\begin{proof}[Proof of \Cref{thm:NG_PROTOCOL_TRAP}]
By \Cref{lem:detailed_balance_path} the path law equals its reversal, so
\[
  \Arr_T(\pi,P)=\DKL{\mathbf{P}^{\pi,P}_T}{\mathbf{P}^{\pi,P}_T}=0.
\]
By \thmNGArrowDPI, $0\le\Arr_T(\pi,P;f)\le\Arr_T(\pi,P)=0$ for every lens $f$.
\end{proof}

\begin{remark}[What the protocol trap does and does not rule out]\label{rem:protocol_trap_reading}
The theorem is about the arrow audit and nothing else.
It allows the observed process $f(X_t)$ to have any amount of temporal structure, and in particular it allows a hidden protocol coordinate to make the observed process pass through recognizable stages.
What it rules out is a preferred direction: under a detailed-balance start, every observed sequence of stages is exactly as likely as the reverse sequence.
The hypothesis concerns the pair $(\pi,P)$ and not $P$ alone, as the next example shows.
\end{remark}

\begin{example}[A reversible kernel with a nonequilibrium start]\label{ex:transient_arrow}
Let $X=\{1,2\}$ and let $P$ have all entries equal to $1/2$.
The uniform law satisfies detailed balance, and started from it the chain has zero arrow.
Started instead from $\mu=\delta_1$, the path $(1,2)$ has probability $1/2$ while its reversal $(2,1)$ has probability $0$, so $\Arr_1(\delta_1,P)=\infty$.
The arrow here records the asymmetry of the initial condition, not a drive: in the decomposition of \Cref{rem:entropy_production}, the total entropy production is $\log 2$ and the boundary term $\DKL{\mu P}{\mu}$ is infinite.
A constant lens hides the arrow completely, while the identity lens shows all of it.
\end{example}

\begin{remark}[Relation to entropy production]\label{rem:entropy_production}
Write $\mu_T=\mu P^T$ for the law at time $T$.
The standard total entropy production over the window compares the forward path law with the reversal of the path law started from $\mu_T$ \citep{Seifert2012,MaesNetocny2003}:
\[
  \Sigma_T(\mu,P)=\DKL{\mathbf{P}^{\mu,P}_T}{\rev_{\#}\mathbf{P}^{\mu_T,P}_T}.
\]
The two quantities are related by
\[
  \Arr_T(\mu,P)=\Sigma_T(\mu,P)+\DKL{\mu_T}{\mu}\qquad\text{in }[0,\infty].
\]
Indeed, the two reversed laws give a path $(x_0,\dots,x_T)$ the masses $\mu(x_T)B$ and $\mu_T(x_T)B$ with the same backward product $B=\prod_tP(x_{t+1},x_t)$.
If all three terms are finite, splitting the logarithm of the likelihood ratio and using that $x_T$ has law $\mu_T$ under the forward path law gives the identity.
Otherwise both sides are infinite: the left side is infinite exactly when some path of positive probability has $\mu(x_T)=0$ or $B=0$; since its endpoint has $\mu_T(x_T)>0$, the first case is exactly the case $\DKL{\mu_T}{\mu}=\infty$ and the second the case $\Sigma_T=\infty$.
Both reversed laws commute with the lens in the same way, so the proof of \thmNGArrowDPI{} also gives
$\DKL{f^{(T)}_{\#}\mathbf{P}^{\mu,P}_T}{\rev_{\#}f^{(T)}_{\#}\mathbf{P}^{\mu_T,P}_T}\le\Sigma_T(\mu,P)$: observation cannot increase the total entropy production either.

If $\pi$ is stationary, the boundary term vanishes, and for $T\ge 1$
\[
  \Arr_T(\pi,P)=\Sigma_T(\pi,P)=T\,\sigma(\pi,P),\qquad
  \sigma(\pi,P)=\sum_{x,x'}\pi(x)P(x,x')\log\frac{\pi(x)P(x,x')}{\pi(x')P(x',x)},
\]
where $\sigma$ is the stationary entropy production rate \citep{Schnakenberg1976,Seifert2012}, the sum runs over pairs with $\pi(x)P(x,x')>0$, and $\sigma=+\infty$ if some such pair has $P(x',x)=0$.
For a stationary chain, \thmNGArrowDPI{} thus bounds the observed arrow by $T\sigma(\pi,P)$.
\end{remark}
